% GENERATED FILE. Edit canonical lessons, then run npm run generate:books.
% canonical-source-hash: fnv1a64:9a1216134899e896
% canonical-lessons: LA-C295-noctu, LA-C295-interdiu, LA-C295-vesperi, LA-C295-mature, LA-C295-dudum

\chapter{Noctū et interdiū}
\label{ch:la-noctu-et-interdiu}
\hlchaptermodality{295}

\begin{chapteropening}
\textbf{By the end of this chapter:} \emph{I can say when in the day things happened: by night, by day, in the evening, early and a while ago.}

Complete the last lesson of chapter 295: I can say when in the day things happened: by night, by day, in the evening, early and a while ago.
\end{chapteropening}

\section[noctū]{noctū — by night, at night}

\label{lesson:LA-C295-noctu}



\begin{quote}
Before the new one: say the Latin for mostly, then the Latin for as often as.
\end{quote}

\subsection*{You'll want to know: noctū}
\textbf{noctū} — \textquotedblleft{}by night, at night\textquotedblright{}.

\begin{quote}
Fūrēs noctū in vīllam intrāvērunt.
\end{quote}

\begin{quote}
\emph{Thieves got into the villa at night.}
\end{quote}

An old case form of \emph{nox} (taught): in the night-time, a word for when thieves and lovers go out.

\begin{cousinweb}
\textbf{noctū} is an old case form of \emph{nox}, \textquotedblleft{}night\textquotedblright{} (EtymDB), which is inherited from Proto-Indo-European \emph{*nók\textsuperscript{w}ts}, \textquotedblleft{}night\textquotedblright{}; English \emph{night} is a cousin.
\end{cousinweb}

\begin{grammarlens}[title={What you've built}]
One more word for telling, explaining and describing at B1.
\end{grammarlens}

\subsection*{Your turn}
\begin{itemize}
\raggedright
  \item \emph{Say it:} \emph{noctū}
  \item \emph{Say it:} \emph{noctū}, once more
  \item \emph{Say it:} a sentence of your own with \emph{noctū}
  \item \emph{Recall:} say the Latin for mostly, then the Latin for as often as, then say \emph{noctū} again
\end{itemize}

\begin{tcolorbox}[breakable,colback=teal!4,colframe=teal!35!black,title={Before you move on}]
What does noctū mean? (\textquotedblleft{}By night, at night\textquotedblright{}.) Say it once more.
\end{tcolorbox}
\section[interdiū]{interdiū — by day, in the daytime}

\label{lesson:LA-C295-interdiu}



\begin{quote}
Before the new one: say the Latin for as often as, then the Latin for by night.
\end{quote}

\subsection*{You'll want to know: interdiū}
\textbf{interdiū} — \textquotedblleft{}by day, in the daytime\textquotedblright{}.

\begin{quote}
Interdiū in agrīs labōrābant, noctū dormiēbant.
\end{quote}

\begin{quote}
\emph{By day they worked in the fields, at night they slept.}
\end{quote}

\emph{inter} + an old form of \emph{diēs}: within the day. It pairs with \emph{noctū}. Caesar and Livy use it; Cicero does not.

\begin{cousinweb}
\textbf{interdiū} joins \emph{inter}, \textquotedblleft{}between, within\textquotedblright{}, and an old form of \emph{diēs}, \textquotedblleft{}day\textquotedblright{} (EtymDB): within the daylight hours.
\end{cousinweb}

\begin{grammarlens}[title={What you've built}]
One more word for telling, explaining and describing at B1.
\end{grammarlens}

\subsection*{Your turn}
\begin{itemize}
\raggedright
  \item \emph{Say it:} \emph{interdiū}
  \item \emph{Say it:} \emph{interdiū}, once more
  \item \emph{Say it:} a sentence of your own with \emph{interdiū}
  \item \emph{Recall:} say the Latin for as often as, then the Latin for by night, then say \emph{interdiū} again
\end{itemize}

\begin{tcolorbox}[breakable,colback=teal!4,colframe=teal!35!black,title={Before you move on}]
What does interdiū mean? (\textquotedblleft{}By day, in the daytime\textquotedblright{}.) Say it once more.
\end{tcolorbox}
\section[vesperī]{vesperī — in the evening}

\label{lesson:LA-C295-vesperi}



\begin{quote}
Before the new one: say the Latin for by night, then the Latin for by day.
\end{quote}

\subsection*{You'll want to know: vesperī}
\textbf{vesperī} — \textquotedblleft{}in the evening\textquotedblright{}.

\begin{quote}
Vesperī omnēs ad cēnam vēnērunt.
\end{quote}

\begin{quote}
\emph{In the evening everyone came to dinner.}
\end{quote}

An old locative of \emph{vesper} (taught): at evening. \emph{Herī vesperī}: yesterday evening.

\begin{cousinweb}
\textbf{vesperī} is an old locative of \emph{vesper}, \textquotedblleft{}evening\textquotedblright{} (Lewis \& Short), used as an adverb.
\end{cousinweb}

\begin{grammarlens}[title={What you've built}]
One more word for telling, explaining and describing at B1.
\end{grammarlens}

\subsection*{Your turn}
\begin{itemize}
\raggedright
  \item \emph{Say it:} \emph{vesperī}
  \item \emph{Say it:} \emph{vesperī}, once more
  \item \emph{Say it:} a sentence of your own with \emph{vesperī}
  \item \emph{Recall:} say the Latin for by night, then the Latin for by day, then say \emph{vesperī} again
\end{itemize}

\begin{tcolorbox}[breakable,colback=teal!4,colframe=teal!35!black,title={Before you move on}]
What does vesperī mean? (\textquotedblleft{}In the evening\textquotedblright{}.) Say it once more.
\end{tcolorbox}
\section[mātūrē]{mātūrē — early; in good time}

\label{lesson:LA-C295-mature}



\begin{quote}
Before the new one: say the Latin for by day, then the Latin for in the evening.
\end{quote}

\subsection*{You'll want to know: mātūrē}
\textbf{mātūrē} — \textquotedblleft{}early; in good time\textquotedblright{}.

\begin{quote}
Agricola mātūrē surrēxit et in agrum īvit.
\end{quote}

\begin{quote}
\emph{The farmer got up early and went to the field.}
\end{quote}

From \emph{mātūrus}, 'ripe, timely'. Its opposite is \emph{sērō}, 'late', the adverb of \emph{sērus} (taught). English \emph{mature}.

\begin{cousinweb}
\textbf{mātūrē} is the adverb of \emph{mātūrus}, \textquotedblleft{}ripe, timely\textquotedblright{} (Lewis \& Short): at the right time, and so early. English \emph{mature} comes from \emph{mātūrus}.
\end{cousinweb}

\begin{grammarlens}[title={What you've built}]
One more word for telling, explaining and describing at B1.
\end{grammarlens}

\subsection*{Your turn}
\begin{itemize}
\raggedright
  \item \emph{Say it:} \emph{mātūrē}
  \item \emph{Say it:} \emph{mātūrē}, once more
  \item \emph{Say it:} a sentence of your own with \emph{mātūrē}
  \item \emph{Recall:} say the Latin for by day, then the Latin for in the evening, then say \emph{mātūrē} again
\end{itemize}

\begin{tcolorbox}[breakable,colback=teal!4,colframe=teal!35!black,title={Before you move on}]
What does mātūrē mean? (\textquotedblleft{}Early; in good time\textquotedblright{}.) Say it once more.
\end{tcolorbox}
\section[dūdum]{dūdum — a little while ago; for a long time}

\label{lesson:LA-C295-dudum}



\begin{quote}
Before the new one: say the Latin for in the evening, then the Latin for early.
\end{quote}

\subsection*{You'll want to know: dūdum}
\textbf{dūdum} — \textquotedblleft{}a little while ago; for a long time\textquotedblright{}.

\begin{quote}
Epistulam, quam dūdum scrīpsī, nunc mittō.
\end{quote}

\begin{quote}
\emph{I am now sending the letter that I wrote a little while ago.}
\end{quote}

With \emph{iam}, 'for a long time now': the track teaches \emph{iamdūdum}.

\begin{cousinweb}
\textbf{dūdum} is built, says Lewis \& Short, on \emph{diū}, \textquotedblleft{}for a long time\textquotedblright{}, and \emph{dum}, \textquotedblleft{}while\textquotedblright{} (EtymDB lists the same parts): it is a while since.
\end{cousinweb}

\begin{grammarlens}[title={What you've built}]
One more word for telling, explaining and describing at B1.
\end{grammarlens}

\subsection*{Your turn}
\begin{itemize}
\raggedright
  \item \emph{Say it:} \emph{dūdum}
  \item \emph{Say it:} \emph{dūdum}, once more
  \item \emph{Say it:} a sentence of your own with \emph{dūdum}
  \item \emph{Recall:} say the Latin for in the evening, then the Latin for early, then say \emph{dūdum} again
\end{itemize}

\begin{tcolorbox}[breakable,colback=teal!4,colframe=teal!35!black,title={Before you move on}]
What does dūdum mean? (\textquotedblleft{}A little while ago; for a long time\textquotedblright{}.) Say it once more.
\end{tcolorbox}
